\clearpage
\section[Abstrak]{}
\sectioncolor{cblue}
\eyebrow{cblue}{Ringkasan}\\[0.3em]
{\bricolagesb\fontsize{19pt}{22pt}\selectfont\color{cink} Abstrak}
\vspace{0.15em}\\
\noindent\textcolor{cblue!35!white}{\rule{\linewidth}{1.1pt}}
\vspace{0.55em}

Laboratorium MGM di Fakultas Ilmu Komputer, Universitas Brawijaya, mempunyai lebih dari 65 anggota yang mempunyai dan mengembangkan berbagai macam proyek. Selama ini, komunikasi dan management hanya mengandalkan spreadsheet atau excel dan pesan WhatsApp maupun Discord. Cara tersebut akan susah scale atau berkembang seiring bertambahnya proyek, dikarenakan kompleksitas yang meningkat. Tim sempat mencoba berlangganan Jira, ClickUp, dan Asana, namun biaya per-pengguna per-bulan tidak wajar untuk skala organisasi yang mempunyai budget minim. Opsi self-hosted seperti Plane juga sudah dicoba, namun tetap terdapat biaya ketika jumlah pengguna melewati batas tertentu, dan tidak menyediakan SSO (SAML/OIDC) dengan paket gratis. Ditambah lagi, hampir semua \textit{project management tools} (PMO) proyek yang dicoba tidak ada yang menyertakan platform komunikasi, sehingga kami tetap harus membayar Slack atau Microsoft Teams secara terpisah untuk berkomunikasi.

\textbf{Atlas} dibuat untuk menyelesaikan persoalan tersebut. Atlas adalah sebuah platform kolaborasi tim yang menyatukan manajemen proyek (PMO, mencakup \textit{kanban board}, \textit{list view}, dan \textit{gantt chart}) dengan sistem komunikasi real-time (RTC), \textit{voice all}, catatan dan whiteboard yang bisa dipakai secara kolaboratif, dalam satu aplikasi yang dapat di-\textit{self-host}, \textit{open source} di bawah lisensi AGPL-3.0, dan tidak punya batasan jumlah pengguna ataupun fitur berbayar. Setiap organisasi cukup menjalankan Atlas di infrastruktur mereka sendiri, mempunyai kendali penuh atas data dari aplikasi, dan tanpa biaya untuk alat manajemen proyek dan alat komunikasi yang terpisah.

\textit{Unique selling points} Atlas dibandingkan hanya prototipe untuk kompetisi HOLOGY adalah statusnya hari ini: aplikasi ini \textbf{sudah berjalan secara produksi} dan digunakan sehari-hari oleh seluruh anggota Laboratorium MGM untuk memanage baik project pribadi maupun project di laboratorium, mulai dari perencanaan, koordinasi tim, hingga komunikasi, bukan sekadar demo yang dibangun untuk perlombaan. Godmode, panel kendali swa-hosting Atlas, memungkinkan admin untu mengatur segala hal yang ada di Atlas (kata sandi, \textit{magic link}, OTP telepon, OAuth, OIDC, SAML), \textit{file storage}, dan tema tampilan tanpa menyentuh kode maupun \textit{env file}.

Atlas mengambil subtema \textbf{Ekonomi Digital} HOLOGY 9.0, Atlas menunjukkan bahwa organisasi Indonesia (laboratorium kampus, komunitas, maupun startup tahap awal) tidak perlu terikat pada langganan perangkat lunak asing yang mahal untuk tetap produktif dan terkoordinasi. Dengan membebaskan biaya operasional yang sebelumnya terserap ke banyak layanan SaaS sekaligus menjaga kedaulatan data di tangan organisasi sendiri, Atlas berkontribusi langsung pada pertumbuhan ekosistem digital yang inklusif, berdaya saing, dan berkelanjutan bagi tim-tim di Indonesia.

\clearpage
